% ============================================================
% ANEXO B — Plan de Entrega Detallado
% ============================================================
\chapter{Plan de Entrega de Artefactos}
\label{anx:plan_entrega}

Este anexo detalla los artefactos a entregar como parte del cierre
formal del Proyecto Demeter, organizados por naturaleza (digital,
f\'isico, conocimiento) y con cronograma asociado.

\section{Artefactos digitales}

\begin{table}[H]
    \centering
    \caption{Artefactos digitales a entregar.}
    \label{tab:digital}
    \begin{tabular}{p{4.5cm}p{8cm}}
        \toprule
        \textbf{Artefacto} & \textbf{Detalle / Ubicaci\'on} \\
        \midrule
        Repositorio p\'ublico & \texttt{github.com/DanielMf31/Proyecto\_Demeter} (privado en el momento de redactar; se abre en la semana 1 del cronograma) \\
        Documentaci\'on T\'ecnica (PDF) & \texttt{docs/Architecture/LaTeX/main.pdf} en el propio repositorio (84 p\'aginas) \\
        Documentaci\'on T\'ecnica (fuentes LaTeX) & Carpeta \texttt{docs/Architecture/LaTeX/} \\
        Informe de Viabilidad (este documento) & Carpeta \texttt{docs/Closure/LaTeX/} \\
        Backups de base de datos & No se entrega \textit{dump}: la base solo contiene datos de prueba y semilla, reproducibles con \texttt{make seed}. El esquema viaja versionado en las migraciones de Alembic \\
        Im\'agenes Docker & Ya publicadas en GitHub Container Registry (\texttt{ghcr.io/danielmf31/proyecto\_demeter/\{backend,frontend\}}); pasan a p\'ublicas junto con el repositorio \\
        Diagramas fuente (.mmd, .drawio) & Carpeta \texttt{docs/diagrams/} en repositorio \\
        \bottomrule
    \end{tabular}
\end{table}

\textbf{Licencias propuestas}:
\begin{itemize}[leftmargin=*]
    \item C\'odigo: \textbf{MIT}. Se descarta Apache 2.0 porque su valor
    diferencial es la cl\'ausula expresa de patentes, irrelevante en un
    proyecto sin cartera de patentes ni intenci\'on de litigio, mientras
    que MIT impone menos fricci\'on a quien quiera reutilizar el c\'odigo
    con fines docentes.
    \item Documentaci\'on: \textbf{CC-BY-SA 4.0} (atribuci\'on +
    compartir igual), de modo que las mejoras sobre el material docente
    reviertan en el mismo r\'egimen abierto.
\end{itemize}

\section{Artefactos f\'isicos}

\begin{table}[H]
    \centering
    \caption{Plan de gesti\'on del hardware f\'isico.}
    \label{tab:fisico}
    \begin{tabular}{p{4cm}p{4cm}p{4.5cm}}
        \toprule
        \textbf{Tipo de hardware} & \textbf{Ubicaci\'on actual} & \textbf{Destino propuesto} \\
        \midrule
        Instalaci\'on hidr\'aulica (4 bombas, 6 electrov\'alvulas, 8 v\'alvulas, 4 barriles de 60\,L, fuente DIN 12\,V) & Almac\'en ESIBot & Puesta a disposici\'on de la ETSIA: financiada con el Proyecto Docente, es patrimonio universitario \\
        Nodos ESP32-S3 con sensores   & Almac\'en ESIBot & Conservaci\'on en la asociaci\'on \\
        Raspberry Pi 4B (gateway)     & Almac\'en ESIBot & Conservaci\'on en la asociaci\'on \\
        Cables, fuentes, accesorios   & Almac\'en ESIBot & Reaprovechamiento en otros proyectos de la asociaci\'on \\
        Carcasas y protecciones IP65  & Almac\'en ESIBot & Reaprovechamiento \\
        Bater\'ias LiPo               & Almac\'en ESIBot & Reaprovechamiento, previa comprobaci\'on de estado de carga \\
        \bottomrule
    \end{tabular}
\end{table}

El detalle por origen de financiaci\'on se recoge en la
tabla~\ref{tab:plan_hardware} del Anexo~\ref{anx:financiero}.

\section{Artefactos de conocimiento e infraestructura}

\begin{table}[H]
    \centering
    \caption{Gesti\'on de cuentas externas e infraestructura.}
    \label{tab:cuentas}
    \begin{tabular}{p{3.5cm}p{4cm}p{5cm}}
        \toprule
        \textbf{Cuenta / Recurso} & \textbf{Estado actual} & \textbf{Acci\'on propuesta} \\
        \midrule
        Cuenta Google Cloud Platform & \textbf{Cerrada}. La VM, el disco persistente y la IP reservada han sido eliminados; la facturaci\'on est\'a detenida & Ninguna. Acci\'on ya ejecutada: era el \'unico gasto recurrente vivo del proyecto \\
        Dominio \texttt{patata.monters.org} & Subdominio del dominio personal del autor, no del proyecto & Retirar el registro DNS. No procede transferencia: el dominio ra\'iz es ajeno al proyecto \\
        Cuenta Cloudflare            & Personal del autor; se usa solo como \textit{tunnel} de acceso al entorno de \textit{staging} & Eliminar el t\'unel de Demeter. La cuenta permanece con el autor \\
        Documentaci\'on interna      & Reside en el propio repositorio (\texttt{docs/}, 14.575 l\'ineas de Markdown), no en una wiki externa & Se entrega con el repositorio; no hay plataforma de terceros que dar de baja \\
        Repositorios privados        & GitHub privado          & Convertir a p\'ublico al cierre \\
        Credenciales y secretos      & Vault privado del autor & Rotar / invalidar antes del cierre p\'ublico \\
        \bottomrule
    \end{tabular}
\end{table}

\begin{warningbox}
Antes de la apertura p\'ublica del repositorio es \textbf{obligatorio}
revisar el hist\'orico Git para asegurar que no contiene credenciales,
claves API, contrase\~nas ni informaci\'on personal sensible.
\end{warningbox}

\subsection*{Resultado de la auditor\'ia de secretos}

La auditor\'ia se ha realizado sobre los 227 \textit{commits} del
repositorio, revisando tanto el \'arbol actual como el hist\'orico
completo. Se resume a continuaci\'on.

\begin{table}[H]
    \centering
    \caption{Hallazgos de la auditor\'ia de secretos del repositorio.}
    \label{tab:auditoria}
    \begin{tabular}{p{1.8cm}p{4.6cm}p{6.2cm}}
        \toprule
        \textbf{Gravedad} & \textbf{Hallazgo} & \textbf{Acci\'on requerida} \\
        \midrule
        \textbf{Cr\'itica} &
        \textit{Token} de t\'unel de Cloudflare v\'alido en
        \texttt{.env.staging}, versionado desde su commit inicial y
        presente tambi\'en en el \'arbol actual. &
        \textbf{Revocar el t\'unel en Cloudflare antes de abrir el
        repositorio.} Una vez revocado, el valor filtrado queda inservible
        y no es necesario reescribir el hist\'orico. \\
        \midrule
        Media &
        Ficheros \texttt{.env} de los cuatro entornos versionados, con
        contrase\~nas de PostgreSQL en claro. &
        Los valores son de desarrollo o marcadores expl\'icitos
        (\texttt{change\_me}), y la base de datos no es accesible desde
        el exterior. Basta con dejar de versionarlos y conservar
        \'unicamente \texttt{.env.example}. \\
        \midrule
        Baja &
        \texttt{Core/auth.py} usa una clave de firma JWT por defecto
        codificada en el fuente si la variable de entorno no est\'a
        definida. &
        No afecta al cierre, pero debe se\~nalarse a quien reutilice el
        c\'odigo: un despliegue que olvide definir la variable firmar\'ia
        sus \textit{tokens} con una clave p\'ublica y conocida. \\
        \bottomrule
    \end{tabular}
\end{table}

\textbf{No se han encontrado} claves de cuenta de servicio de Google
Cloud, claves privadas SSH o TLS, ni datos personales de terceros en
ning\'un punto del hist\'orico.

Con la revocaci\'on del \textit{token} de Cloudflare --- \'unica acci\'on
bloqueante --- el repositorio puede abrirse sin reescribir el hist\'orico,
lo que evita invalidar los identificadores de \textit{commit} citados a
lo largo de esta documentaci\'on.

\section{Cronograma detallado}

\begin{table}[H]
    \centering
    \caption{Cronograma de entrega tras la aceptaci\'on del cierre.}
    \label{tab:cronograma_detalle}
    \begin{tabular}{p{2cm}p{11cm}}
        \toprule
        \textbf{Semana} & \textbf{Acciones} \\
        \midrule
        1 & Auditor\'ia de seguridad del repositorio (eliminaci\'on de
            secretos del hist\'orico). Selecci\'on de licencias.
            Apertura del repositorio en GitHub p\'ublico. Anuncio
            interno en la asociaci\'on. \\
        \midrule
        2 & Entrega formal de la documentaci\'on a la Universidad
            (PDFs + fuentes). Env\'io del informe de cierre completo.
            Reuni\'on de cierre con el Departamento de Agronom\'ia de la
            ETSIA, promotor del Proyecto Docente, y con la junta directiva
            de ESIBot. \\
        \midrule
        3 & Inventario f\'isico definitivo del hardware. Devoluci\'on /
            cesi\'on seg\'un origen de financiaci\'on. Firma del acta
            de entrega de hardware. \\
        \midrule
        4 & Baja del registro DNS y del t\'unel de Cloudflare. Acta de
            cierre t\'ecnico. El cierre de la infraestructura cloud
            (eliminaci\'on de la VM, el disco y la IP reservada) ya est\'a
            ejecutado, por lo que esta semana solo recoge la
            formalizaci\'on documental. \\
        \bottomrule
    \end{tabular}
\end{table}

\section{Documento de cierre formal}

Como conclusi\'on del proceso, se solicita la firma de un \textbf{acta
de cierre} entre la Universidad de Sevilla y la Asociaci\'on
Universitaria ESIBot, que recoja al menos:

\begin{itemize}[leftmargin=*]
    \item Reconocimiento de la conclusi\'on del Proyecto Docente Demeter.
    \item Listado de artefactos entregados.
    \item Disposici\'on del hardware adquirido con fondos institucionales.
    \item Reconocimiento del valor docente generado.
    \item Disponibilidad futura de los autores para consultas puntuales.
    \item Apertura a futuras colaboraciones bajo nuevos formatos.
\end{itemize}

Este acta constituye el cierre formal del proyecto y referencia
institucional de la colaboraci\'on, facilitando que ambas partes
mantengan relaci\'on operativa para futuras iniciativas.
